\documentclass[10pt,twocolumn]{article}
\usepackage[T1]{fontenc}
\usepackage[utf8]{inputenc}
\usepackage[letterpaper,margin=18mm,columnsep=7mm]{geometry}
\emergencystretch=2em
\usepackage{lmodern,microtype,amsmath,amssymb,graphicx,booktabs,tabularx,array,xcolor,url}
\usepackage[numbers,sort&compress]{natbib}
\usepackage[hidelinks]{hyperref}
\usepackage{stfloats}
\graphicspath{{figures/}}
\newcommand{\ua}{u_A}
\newcommand{\am}{a_m}
\newcommand{\om}{\Omega_8}
\hypersetup{pdftitle={Resource-consistent shape optimization of formed corrugated media: analytic curvature and finite-scenario qualification},pdfauthor={Ricardo Fitas; Sergio Manuel Oliveira Tavares; Oliver Weeger}}
\title{Resource-consistent shape optimization of formed corrugated media: analytic curvature and finite-scenario qualification}
\author{Ricardo Fitas$^{1,2}$, S\'ergio Manuel Oliveira Tavares$^1$, Oliver Weeger$^2$\\[3pt]
\small $^1$Centre for Mechanical Technology and Automation (TEMA), Department of Mechanical Engineering,\\
\small University of Aveiro, Campus Universit\'ario de Santiago, 3810--193 Aveiro, Portugal\\
\small $^2$Cyber-Physical Simulation, Department of Mechanical Engineering, Technical University of Darmstadt,\\
\small Dolivostrasse 15, 64293 Darmstadt, Germany\\
\small Correspondence: \texttt{rfitas99@gmail.com}}
\date{Version 2 research manuscript, 4 September 2026}
\begin{document}
\raggedbottom
\maketitle
\begin{abstract}
Shape optimization of corrugated media becomes physically interpretable when material, initial envelope and geometric admissibility are compared consistently. We reformulate the design resource as actual medium volume per initial projected area and evaluate NURBS curvature at stationary points and both sides of its knots. An audit of 572 archived finite-element geometries finds 16 radius-screen violations previously admitted by smoothed curvature. Correcting only the material coordinate reduces an archived 31-member material--stress front to one member under its historical eligibility conditions. A new experiment performs 449 direct compression paths for sinusoidal, Fourier and NURBS shapes at fixed pitch, height and medium resource. Three paired 64-call comparisons show no consistent NSGA-II advantage over Sobol sampling. Ten finalists receive 270 additional paths across 27 geometric and material sensitivity scenarios. Four remain qualified; joint perturbations reject one candidate that passed all 11 axial scenarios. At 24 elements, the qualified NURBS energy extreme retains 7.91\% greater minimum stored potential per footprint than the sinusoid, with 3.38\% greater maximum stress utilization in the sampled set. The nominal minimum-stress design and normalized knee fail geometric qualification. These results establish a practical selection principle: resource matching, knot-sided curvature and scenario enrichment materially change which nominal profile trade-offs remain admissible. The findings are conditional numerical design evidence for an attached medium between rigid surfaces; stored potential is not irreversible crush dissipation.

\end{abstract}
\noindent\textbf{Keywords:} corrugated medium; equal material resources; rational B-splines; curvature; nonlinear compression; numerical robustness.

\section{Introduction}
A useful shape optimum must compare the same physical resources and survive the numerical checks that define its admissibility. These requirements are particularly consequential for corrugated paper, where pitch, height, thickness, contact and local curvature influence different response measures. Detailed corrugated-board models and experiments already establish the importance of geometry and boundary conditions \citep{Talbi2009,Park2021Flat,DiRusso2023}. Rational B-splines and evolutionary multiobjective search provide established foundations for exploring this design space \citep{Piegl1997,Deb2002}. Their engineering value depends on the physical comparability and geometric admissibility of the resulting candidates.

This study asks a more specific question: which shape advantages remain after fixing the initial envelope and actual medium volume per projected area, evaluating curvature at its analytic extrema, and perturbing the manufactured geometry and material parameters? The answer is developed through an audit of an existing 572-geometry finite-element ledger and a new, directly evaluated shape comparison. The two components have different evidential roles. The audit diagnoses consequences of definitions and numerical constraints in the archived study. The new experiment compares sinusoidal, two-parameter Fourier and six-parameter NURBS shapes under the same nominal physical resources and finite-element protocol.

Building on sensitivity-aware multiobjective optimization \citep{DebGupta2006}, the study establishes three concrete contributions: a resource correction that changes the archived material--stress selection, an analytic curvature audit that identifies previously admitted radius violations, and an equal-budget direct-FEM comparison followed by finite-scenario qualification. Together they identify which apparent shape advantages survive physically comparable accounting and tighter geometric evaluation. Complete call ledgers connect these conclusions to the executed calculations.

\section{What is being optimized?}
\subsection{Medium resource and stored potential}
Let $P$ be initial pitch, $H$ initial centreline height, $b$ specimen width, $t$ medium thickness and $L_s$ the centreline arc length of one period. The actual medium volume divided by its initial projected area is
\begin{equation}
 \am=\frac{btL_s}{bP}=\frac{tL_s}{P},\qquad
 \frac{m_m}{A_0}=\rho_m\am,\quad A_0=bP. \label{eq:resource}
\end{equation}
The second equality assumes a common medium density $\rho_m$; density is not measured or assigned in this numerical study. We consequently report $\am$ in mm, without converting it into an asserted grammage.

If two liners of the same thickness and density as the medium are included in bookkeeping, the corresponding volume per initial footprint is
\begin{equation}
 a_b=t(L_s/P+2)=H\phi_b,\qquad
 \phi_b=\frac{t(L_s+2P)}{PH}. \label{eq:board}
\end{equation}
Thus $\phi_b$ is an envelope volume fraction. It is not proportional to areal mass when $H$ varies. For independently specified liners, $m_b/A_0=\rho_m tL_s/P+\rho_b t_b+\rho_t t_t$. In the new experiment, only the medium is mechanically represented, and Eq.~\eqref{eq:resource} is the resource constraint. Fixed liner contributions could subsequently be added to its mass accounting; they are not part of the computed deformation energy.

The objective pair uses terminal stored potential per initial footprint and a stress norm:
\begin{equation}
 \max\ \ua=\frac{U(0.2H)}{bP},\qquad
 \min\ \om=\frac{1}{\sigma_y}
 \left(\frac{\sum_e L_e\sigma_e^8}{\sum_e L_e}\right)^{1/8}. \label{eq:objectives}
\end{equation}
Here $\sigma_e$ is the largest absolute longitudinal extreme-fibre stress in element $e$. The ordinate $\ua$ is displayed in J\,m$^{-2}$: 1 N\,mm$^{-1}=1000$ J\,m$^{-2}$. It measures recoverable potential in the specified smooth deformation law. It does not measure irreversible crushing energy, nor is it interchangeable with experimental loading work.

\subsection{Audit of the archived interpretation}
All latest supplied manuscript files are treated as the v1 baseline and retained unchanged. The repository's mesh-24 ledger contains 572 successful geometries and 732 case classifications; its 58, 31 and 121 archived front memberships are reproduced for C037--C039. These counts describe fronts within visited candidate sets.

Table~\ref{tab:resource_audit} recomputes the two comparisons most affected by Eq.~\eqref{eq:board}. S4 has 27.43\% lower $\phi_b$ than S6 but 12.71\% greater medium volume per footprint. Its larger initial height also changes absolute platen travel at the common engineering strain. The supposed small material change from S7 to S9 corresponds to a 59.69\% reduction in medium volume per footprint and a 95.47\% reduction in stored potential per footprint. Neither comparison isolates a shape advantage at equal resources.

\begin{table}[t]
\centering\small
\caption{Archived response ratios, expressed as $100(A/B-1)$. The same-thickness liner column is hypothetical bookkeeping.}\label{tab:resource_audit}
\begin{tabular}{lrr}
\toprule Quantity & S4/S6 [\%] & S9/S7 [\%]\\\midrule
Envelope fraction $\phi_b$ & $-27.43$ & $+1.48$\\
Initial height $H$ & $+42.02$ & $-45.00$\\
Medium resource $\am$ & $+12.71$ & $-59.69$\\
Same-$t$ board resource $a_b$ & $+3.07$ & $-44.19$\\
Stored potential/area $\ua$ & $+40.45$ & $-95.47$\\
Normalized potential $\Psi_U$ & $+24.61$ & $-88.75$\\
Stress utilization $\om$ & $+10.52$ & $-52.75$\\\bottomrule
\end{tabular}
\end{table}

A particularly clear diagnostic is obtained without changing the 129 historically eligible C038 rows or their original normalized-potential floor. Replacing only $\phi_b$ by $\am$ in the material--stress objective pair reduces the 31-member front to S6 alone: S6 dominates every other eligible row in these two coordinates. This is a conditional reclassification, not a fixed-height optimum, because the historical eligibility conditions still mix heights. It motivates the resource-matched experiment below.

\begin{figure*}[t]
\centering\includegraphics[width=.98\textwidth]{archive_resource_audit.pdf}
\caption{Why the material definition changes selection. Left: the historically eligible C038 candidates in the original fraction--stress coordinates. Right: the identical candidates with actual medium volume per initial footprint. The original 31-member front collapses to S6 under the corrected material coordinate. Historical eligibility is held fixed in this diagnostic.}\label{fig:archive}
\end{figure*}

\section{Analytic geometric admissibility}
\subsection{A knot-sided radius calculation}
The archived implementation estimates curvature from smoothed finite differences of a sampled NURBS curve. The v1 main text described analytic derivatives, while its appendix and code used smoothing. The present implementation makes that distinction explicit and leaves the archived evaluator unchanged.

A degree-two rational span can be written $\mathbf c=\mathbf N/w$, with quadratic homogeneous coordinate polynomials. Using a local span coordinate, define
\begin{align}
 \mathbf A&=\mathbf N'w-\mathbf Nw',&D&=\mathbf A\cdot\mathbf A,\nonumber\\
 K&=w^2\det(\mathbf A,\mathbf A').
\end{align}
Its signed curvature is $K/D^{3/2}$. Candidate interior extrema satisfy
\begin{equation}
 2K'D-3KD'=0. \label{eq:curvature_extrema}
\end{equation}
The implementation evaluates every real polynomial root in each span, both span endpoints, and both sides of each internal knot. Degree-two curves are $C^1$ at simple knots; their second derivatives need not agree there. The periodic seam is included. The maximum absolute curvature over these candidates gives $R_{\min}^{-1}$. This is a floating-point stationary-point calculation, not an interval-arithmetic certificate.

For Fourier shapes and smooth perturbations of NURBS, analytic first and second derivatives are also used, but local maxima of curvature are bracketed on 257 points per smooth span and refined by bounded scalar minimization. The supplementary verification distinguishes this numerical extremum search from the polynomial calculation.

\subsection{Consequences for archived feasibility}
Of the 572 archived geometries passing the smoothed $R_{\min}\geq0.9$ mm screen, 16 fail the analytic screen. The largest smoothed-radius overestimate is 7.90\%, although the median is only 0.0803\%. S1 and S3, both archived objective extremes, move below 0.9 mm. A small typical error is therefore insufficient assurance near an active inequality.

The analytic coordinates agree with the archived dense NURBS construction to $8.9\times10^{-15}$ mm in 40 targeted checks, comprising all 16 changed classifications and 24 sampled archived rows. Analytic tangents agree with centred finite differences to relative error $1.2\times10^{-10}$ away from knots. A much denser radius evaluation, at 8193 points per span, exceeds the stationary-point minimum-radius estimate by at most relative $2.7\times10^{-8}$ in these checks. The common spacing rescaling and periodic tangent identities also pass. These are numerical verification results, not material or manufacturing validation.

\begin{figure*}[t]
\centering\includegraphics[width=.98\textwidth]{curvature_audit.pdf}
\caption{Radius-screen audit of 572 archived designs. The 0.9 mm criterion is applied without display rounding. The analytic method evaluates knot sides and stationary points; the historical method uses smoothed finite differences. The selected-design panel shows why near-threshold representatives require explicit rechecking.}\label{fig:radius}
\end{figure*}

The inherited 0.9 mm forming-radius threshold is retained as an illustrative design constraint. Passing it does not establish process capability. In particular, the medium is assumed stress-free after forming; residual forming stress, springback and damage are outside the mechanics below.

\section{New fixed-resource experiment}
\subsection{Shape families and resources}
Every nominal design has $P=7.9$ mm, $H=3.0$ mm and $\am=0.27$ mm. Thickness is solved from $t=\am P/L_s$ and must lie in $[0.15,0.25]$ mm. Thus shape-dependent arc length changes required thickness, rather than silently changing material quantity. Equal resources refer to the initial configuration; the compression boundary model allows the horizontal pitch to evolve.

The sinusoidal reference is $y=H[1+\cos(2\pi x/P)]/2$, whose radius is $P^2/(2\pi^2 H)=1.053909$ mm. The two-parameter smooth comparator is
\begin{align}
 y=\frac H2[1+\cos\theta
 &+a(\cos2\theta-1)\nonumber\\
 &+b(\cos3\theta-\cos\theta)],\quad\theta=2\pi x/P, \label{eq:fourier}
\end{align}
with $a\in[-0.12,0.12]$ and $b\in[-0.08,0.08]$. The crown and trough heights remain exactly $H$ and zero; monotonicity on each half-period is checked.

The NURBS family retains the v1 ten-control-point periodic quadratic construction. Its five positive spacings are normalized by their sum, and two rational-weight parameters lie in $[0,1]$. The new search removes the redundant common spacing scale: $x_j=\exp[4(z_j-1/2)]$ for $j=1,\ldots,4$, $x_5=1$, and the remaining two coordinates are weight parameters, with $\mathbf z\in[0,1]^6$. This declares a six-dimensional subset of the historical ratio domain, rather than claiming to repeat its entire ten-variable search. The complete construction is given in the supplement and executable source.

\subsection{Mechanics and boundary conditions}
A preformed medium of width $b=25$ mm is discretized by corotational assumed-strain Timoshenko beam elements \citep{Crisfield1991}. Constant axial strain and curvature are integrated through thickness using five Gauss points, and the mean shear strain is integrated once per element. The longitudinal deformation-potential density is
\begin{equation}
 W(e)=\frac{hE e^2}{2}+(1-h)\frac{\sigma_y^2}{E}
 \left[\sqrt{1+(Ee/\sigma_y)^2}-1\right], \label{eq:energy}
\end{equation}
with $E=2899$ MPa, $G=E/55$, $\sigma_y=60$ MPa and $h=0.02$. The stress is $W'(e)$. This smooth, single-valued law has no committed plastic strain and no dissipative unloading loop. The terminal $U$ in Eq.~\eqref{eq:objectives} is exactly the potential differentiated by the solver.

The exact trough is persistently attached vertically to the lower rigid surface. Both periodic crowns are persistently attached vertically to the upper rigid surface; their vertical displacement and rotations are tied. One crown's horizontal displacement fixes rigid translation, while the opposite crown's horizontal displacement is free. The remaining nodes obey unilateral vertical bounds \citep{Wriggers2006}; tangential slip is free. Surface offsets are $t/2$. These are attached rigid-liner boundary idealizations, not fully detachable contact or a flexible two-liner board. Self-contact, friction, debonding and material history are absent.

Crown and trough positions are retained as nodes at every mesh, and the remaining elements are allocated by the two arc lengths. Compression is prescribed at $\delta/H=0,0.05,0.10,0.15,0.20$. The inherited continuation uses L-BFGS-B with the original constrained fallback and step subdivision checks. A state requires optimizer success and normalized projected gradient at most $5\times10^{-4}$, using force scale $\sigma_ybt$. Every path must have all accepted nominal states. A converged minimization path is not a proof of global equilibrium selection, complete branch discovery or positive reduced-Hessian inertia.

\subsection{Budgets and comparison rules}
The NURBS experiment compares direct NSGA-II with scrambled Sobol sampling \citep{Deb2002,JoeKuo2008} for seeds 4101--4103. Each method receives exactly 64 expensive, 16-element compression paths per seed. Their first 16 feasible Sobol designs are identical for each paired seed. NSGA-II then uses rank/crowding tournaments, simulated binary crossover with distribution index 15, polynomial mutation with index 20 and probability $1/6$, and elitist rank/crowding selection, over four 16-member batches in total. Geometrically rejected proposals are screened before FEM and counted separately. Failed FEM paths consume their allotted call and are retained as unsuccessful outcomes.

A separate 64-call Sobol run covers the Fourier family, and one call evaluates the sine. The complete nominal budget is therefore 449 direct paths. There is no hidden surrogate-training charge in this experiment. The comparator addresses one family and one small expensive-evaluation budget, with the two complete direct-search procedures assessed under identical mechanics.

Hypervolume uses $(u_{A,\mathrm{sine},16}/\ua,\om)$ and fixed reference $(2,1)$, where the sine scale was obtained in the timing pilot. Results are pooled over all evaluated feasible candidates, including designs discarded from a final population. Histories use expensive path calls as the horizontal coordinate. Individual paired outcomes and elapsed times are reported, with no significance or equivalence claim from three pairs. The protocol was written before the comparative solves, after a disclosed timing/geometry pilot; it was not externally preregistered.

\subsection{Finite-scenario qualification}
The shortlist is selected by a fixed rule: the sinusoid; each family's energy extreme, stress extreme and normalized knee; the largest-radius NURBS candidate; and four evenly spaced members of the pooled nominal front. Identical family/parameter vectors are deduplicated. The same shortlist is first evaluated in 11 scenarios (nominal plus signed single-axis perturbations) and then enriched by 16 joint Sobol perturbations generated with seed 4151.

The sensitivity ranges are $\pm3\%$ in height, $\pm5\%$ in thickness, $\pm10\%$ in $E$ and $G$ together, $\pm10\%$ in $\sigma_y$, and a smooth forming perturbation $\eta\in[-0.02,0.02]$:
\begin{equation}
 y_\eta(u)=y(u)+\eta H\sin^2(2\pi q),\qquad
 q=\frac{u-u_0}{u_1-u_0}. \label{eq:perturbation}
\end{equation}
This perturbation preserves the reference crown and trough coordinates and periodic tangent closure; half-period monotonicity is rechecked. Thickness is not renormalized after perturbation: manufacturing variation is allowed to change the actual material resource, which is recorded for every scenario. Each perturbed specimen reaches the same engineering strain, 0.20. These illustrative ranges are not inferred probability distributions or calibrated process tolerances.

Each shortlisted design is re-solved directly at 24 elements in every scenario. Define
\begin{equation}
\begin{aligned}
 \underline{u}_{A,\mathcal S}&=\min_{s\in\mathcal S}\ua(s),\\
 \overline{\Omega}_{8,\mathcal S}&=\max_{s\in\mathcal S}\om(s),\\
 \underline R_{\mathcal S}&=\min_{s\in\mathcal S}R_{\min}(s).
\end{aligned}\label{eq:robust}
\end{equation}
Finite-set qualification requires every path to be accepted, every geometry screen to pass and $\underline R_{\mathcal S}\geq0.9$ mm. These extrema define the comparison within the shortlist and sampled scenarios. Enrichment can worsen an extremum or remove feasibility, providing an explicit test of whether a design decision remains stable when the sensitivity set grows.

\section{Results of the new computations}
\subsection{Equal-resource shape trade-offs}
All 449 nominal solver paths meet the inherited acceptance checks. Nominal thickness varies from 0.18919 to 0.20955 mm because arc length is compensated to maintain $\am=0.27$ mm. Figure~\ref{fig:new_front} shows a continuous energy--stress trade-off across the visited shapes. The 16-element sinusoidal reference has $\ua=12.8137$ J\,m$^{-2}$ and $\om=0.51505$. The NURBS energy extreme V01 increases stored potential by 7.79\% with 3.91\% greater stress utilization. The NURBS stress extreme V02 reduces utilization by 10.48\%, at the cost of 22.92\% less potential. These are physically matched shape trade-offs, rather than savings obtained by varying height or material quantity.

The Fourier family provides a useful simpler comparator. Its energy extreme exceeds the sine by only 0.66\%, with 0.09\% greater stress utilization; its stress extreme lowers utilization by 2.49\% while sacrificing 5.81\% of potential. Increased parameter count mainly enlarges the observed NURBS trade-off range. It does not produce a shape that improves both objectives over every simpler reference.

\begin{figure*}[t]
\centering\includegraphics[width=.98\textwidth]{fixed_resource_fronts.pdf}
\caption{Fixed-resource direct-FEM comparison at mesh 16. Left: all accepted nominal evaluations and each family's pooled candidate front; the sinusoid is an explicit simpler reference. Right: median and full seed range of pooled-candidate hypervolume. Each paired algorithm run receives exactly 64 expensive paths and the same initial 16 candidates.}\label{fig:new_front}
\end{figure*}

The paired final hypervolume changes for NSGA-II relative to Sobol are $+1.10\%$, $-4.64\%$ and $-1.06\%$ for seeds 4101--4103. Mean values are 0.53906 and 0.54759, respectively. Thus the tested evolutionary search has no consistent front-quality advantage at this budget. It does require fewer geometric proposals: 1571 in total across its three runs, compared with 6091 for Sobol. Both consume 192 expensive paths. Table~S1 gives each run's counts, indicator and elapsed time; differences in geometric filtering and solver conditioning are included in the timings. The result supports retaining a direct sampling baseline when assessing a more elaborate design pipeline.

\subsection{Which nominal choices survive perturbations?}
The ten shortlisted designs receive all 270 prescribed scenario paths. Every submitted solver path meets the inherited algorithmic acceptance checks, while 39 scenario geometries fail the declared admissibility tests: 32 violate the 0.9 mm radius and seven lose half-period monotonicity. The latter all involve Fourier stress extreme V05, despite its minimum radius remaining above 1.11 mm. The largest reference-envelope overshoot in these rejected geometries is 0.00133 mm. This rejection concerns the declared single-trough shape family, rather than an experimentally measured forming failure.

Five candidates pass all 11 axial scenarios. Joint enrichment removes V09, whose minimum radius decreases from 0.90676 to 0.89750 mm. Four designs remain qualified in all 27 scenarios: the sinusoid V00 and NURBS candidates V01, V07 and V08. The nominal minimum-stress choice V02 and nominal normalized knee V03 are rejected, with scenario minimum radii 0.87860 and 0.88926 mm. A knee chosen only in the nominal objective plane therefore does not automatically provide a defensible design margin.

\begin{figure*}[t]
\centering\includegraphics[width=.98\textwidth]{scenario_qualification.pdf}
\caption{Finite-scenario qualification. Left: nominal mesh-24 responses move to the minimum-potential/maximum-stress coordinates over 27 scenarios; green identifies complete geometric qualification. The two extrema may occur in different scenarios. Right: minimum radius after axial and joint enrichment. V05 also violates the separately checked monotonicity condition.}\label{fig:scenario}
\end{figure*}

All four qualified candidates are nondominated within this finite shortlist. Table~\ref{tab:qualified_main} gives their sampled extrema. V01 retains 7.91\% greater minimum potential than the sine, with 3.38\% greater maximum stress utilization. V08 offers 2.97\% lower maximum utilization, with 12.20\% less minimum potential. V07 trades response for a substantially larger radius margin. The result is a qualified menu of engineering choices, rather than one preference-independent optimum.

\begin{table*}[t]
\centering\small
\caption{Qualified designs over the enriched 27-scenario set, evaluated at mesh 24. Material variation is part of the sensitivity study; nominal resource is matched.}\label{tab:qualified_main}
\begin{tabular}{llrrrr}
\toprule ID & Role & Minimum $R$ [mm] & Minimum $\ua$ [J\,m$^{-2}$] & Maximum $\om$ & Nominal $\ua$ [J\,m$^{-2}$]\\\midrule
V00 & Sinusoidal reference & 0.96675 & 10.35180 & 0.58173 & 12.69798\\
V01 & NURBS energy extreme & 0.93598 & 11.17054 & 0.60141 & 13.68619\\
V07 & Largest-radius NURBS & 1.57747 & 9.30197 & 0.57172 & 11.42222\\
V08 & Qualified low-stress alternative & 0.99724 & 9.08890 & 0.56443 & 11.15943\\\bottomrule
\end{tabular}
\end{table*}

The varied medium resource ranges from 0.25405 to 0.28350 mm across the scenarios. This is recorded explicitly because re-adjusting thickness to restore nominal material after a manufacturing perturbation would define a different uncertainty problem. Minimum potential and maximum stress generally occur in different scenarios. Their pairing in Eq.~\eqref{eq:robust} deliberately represents two finite-set extrema, not one physically realized state.

\begin{figure*}[t]
\centering\includegraphics[width=.98\textwidth]{qualified_profiles.pdf}
\caption{Qualified examples under the same nominal initial resource and envelope. Dashed lines show reference geometry; solid lines show the mesh-24 deformed centreline colored by longitudinal extreme-fibre stress. Reaction plots show nominal response and the full 27-scenario range. V07 is the normalized finite-scenario knee within the qualified NURBS shortlist.}\label{fig:profiles}
\end{figure*}

\subsection{Numerical resolution and full evaluation cost}
The nominal search, scenario qualification, targeted refinement and fine-mesh retries together contain 764 expensive path attempts, in addition to three disclosed pilot calls. All 449 nominal and 270 scenario solver paths satisfy the inherited algorithmic checks. The initial targeted refinement and its first extension supply 21 mesh-32 paths, six mesh-40 paths and six nine-point load paths. Maximum 24--32 changes are 1.20\% in potential, 0.334\% in stress utilization and 6.35\% in reaction. The reaction sensitivity of V01 motivated the additional fine checks.

All six default mesh-64 paths reached the inherited iteration limits without satisfying the full success predicate. They remain recorded as rejected. Six retries increased the primary iteration cap from 1800 to 6000, leaving the mechanics, load protocol, tolerances and success predicate unchanged; all six then passed. No rejected high-resolution response enters the comparisons.

The sine changes by at most 0.171\% in potential between meshes 40 and 64. For V01, potential changes by 1.83--1.85\%, stress utilization by 0.65--0.89\%, and reaction by 0.58--0.77\% in absolute value over this interval. At the three checked mesh-64 scenarios, its potential advantage over the matched sine remains positive, 13.11--14.33\%, with 4.87--6.07\% greater stress utilization. These are checks at the original nominal and extremal scenarios, not a repeated 27-scenario scan at mesh 64. The sign of the interpreted trade-off persists, while its magnitude remains resolution-dependent. The 7.91\% and 3.38\% results therefore specifically describe the mesh-24 finite-scenario calculation, rather than converged continuum values.

Increasing the nominal load path from five to nine points changes terminal potential by at most 0.000098\%, stress utilization by 0.0253\% and reaction by 0.290\% in the six checks. Continuous resource matching is exact by construction; the largest polygonal-mesh resource deficit is 0.758\% at mesh 16 and 0.335\% at mesh 24. Thus small differences near the nominal front must be read alongside mesh sensitivity. The finer evidence, failed attempts and state arrays are retained in the supplement and numerical package.



\section{Interpretation across the four studies}
This design study occupies a different evidential role from its three companion manuscripts. Paper B \citep{CompanionB} supplies specimen-level compression-work and sequential load-loss evidence. Its common observation windows motivate explicitly declaring the compression endpoint here, but experimental loading work is not the deformation potential in Eq.~\eqref{eq:energy}; paper category is also not an identified probability law for our perturbations. Shared original specimens are reused experimental evidence, not independent validations of several models.

Paper C \citep{CompanionC} examines state representation, exact reduced mechanics and the full cost of qualified computational responses. Its results motivate charging training, verification and state recovery whenever a surrogate is used for design. The present new comparison instead charges direct FEM paths. Any later acceleration should reproduce the present resource definitions and solver-acceptance outcomes under these boundary conditions before its speedup is transferred.

Paper C2 \citep{CompanionCtwo} demonstrates why close reactions can coexist with different configurations and compares finite searches for certified elastica states. Its ten-wave fixed-span elastica, support ties and constitutive history differ from the single-period model here. It is a warning against using force agreement as a state certificate, not evidence that our scenario extrema contain all branches. Our qualification varies geometry and material along the selected continuation protocol; it does not establish robustness to alternative equilibria.

The four papers thus connect through observability, resources and explicitly bounded numerical evidence. Measurements define what occurred, mechanics define what a state predicts, qualification defines when a reduced calculation is credible, and optimization selects among physically comparable candidates. The interfaces require matching specimens, boundary conditions and state variables rather than transferring a validation statistic.

\section{Limitations and next physical test}
The new experiment fixes one initial pitch, height and medium resource; it covers two bounded smooth shape families and one loading horizon. The NURBS log-ratio bounds are explicit and do not exhaust all periodic profiles. The candidate-level FEM comparison and finite-scenario front remain conditional on these choices. The original linear homogenized modulus and geometric-inertia studies are retained in the v1 archive, but they are not substituted for nonlinear response in the new comparison.

Mesh and load-step checks qualify selected responses, not every continuous optimum. The radius calculation is verified numerically; the 0.9 mm constraint and sensitivity ranges have no measured process-capability interpretation. The inherited smooth potential excludes damage, friction, flexible liners, adhesive constitutive behavior, residual forming stress and path-dependent plasticity. Different stable configurations may also be reached under different histories or disturbances. Consequently, numerical advantages do not yet establish superiority of a commercial board construction.

A discriminating next experiment would fabricate the sinusoidal reference, a nominally attractive candidate and a scenario-qualified candidate at measured common medium mass and initial envelope. Formed geometry, actual thickness, synchronized side images, transmitted force and loading/unloading work should be recorded on a common displacement protocol. Such measurements would test resource matching, forming feasibility and configuration selection separately. They cannot be replaced by the number of optimizer runs or the size of an archived front.

\section{Conclusions}
A physically matched design comparison changes both the interpretation and selection of corrugated profiles. In the archived nonlinear campaign, actual medium volume per footprint reverses an asserted material saving and reduces the historically eligible C038 material--stress front from 31 members to S6 alone. Analytic knot-sided curvature rejects 16 of the 572 previously admitted geometries. These effects are consequences of resource definitions and geometric evaluation, rather than optimizer choice.

The new direct-FEM experiment fixes pitch, height and medium volume per initial area, then compares simple smooth references with NURBS. At the tested 64-path budget, three paired runs provide no consistent NSGA-II advantage over Sobol. The richer family expands the observed energy--stress trade-off, while the sinusoid remains a competitive reference.

Finite-scenario qualification changes the shortlist: four of ten candidates remain admissible after 27 perturbations, and enrichment removes a candidate that passed all axial tests. At 24 elements, a qualified energy-focused NURBS profile retains a 7.91\% sampled minimum-potential advantage over the sine with a 3.38\% sampled maximum-stress cost. Finer checks preserve the sign of this energy--stress trade-off in the inspected states, while its magnitude remains resolution-dependent. The nominal minimum-stress design and knee do not survive the geometric checks. Explicit resources, analytic geometric screening and enriched perturbation tests therefore provide the evidence needed to choose which profiles merit controlled physical prototyping.


\section*{Data and code availability}
The v1 manuscript, figures and bibliography supplied for this revision are preserved with file hashes. The public code baseline is commit \texttt{9144fee3} of \url{https://github.com/rfitas-lab/corrugated-board-size-shape-optimization}. The accompanying v2 package contains the analytic geometry evaluator, inherited-mechanics adapter, frozen protocol, every expensive call and geometric-screen count, solver failures, saved states, scenario and refinement ledgers, figure generators and manuscript sources. All v2 results are separate from the archived campaign. The supplement provides commands and provenance; availability of an eventual public v2 revision is distinct from the baseline repository link.

\section*{Author declarations}
\textbf{Funding:} This work was supported by a University of Aveiro research grant (BI/PRR/13036/2026) within Agenda ATE---Alliance for the Energy Transition, project C644914747-00000023, funded by the Portuguese Recovery and Resilience Plan (PRR) and European Funds NextGenerationEU through the Agendas for Business Innovation programme.
\textbf{Competing interests:} The authors declare no competing interests.
\textbf{Author contributions:} R.F.: conceptualization, methodology, software, formal analysis, data curation, visualization and writing---original draft. S.M.O.T.: supervision, validation and writing---review and editing. O.W.: methodology, supervision, validation and writing---review and editing.
\bibliographystyle{unsrtnat}
\bibliography{references}
\end{document}
